\newcommand{\F}{\mathbb{F}}
\newcommand{\C}{\mathbb{C}}
\newcommand{\an}{\mathscr{A}} 
\newcommand{\Q}{\mathbb{Q}}
\renewcommand{\hom}{\operatorname{Hom}}
\newcommand{\ehom}{\operatorname{End}}
\colorlet{gris}{black!70}

\-\vspace{-0.3cm}
  {\hrule \vspace{-0.15cm}
  \centering 2025-II \vspace{0.1cm}
  \hrule}

1. Defina ou enuncie em detalhe cada um dos ítens em baixo: \\
(a) Aplicação diff num ponto \(\x\). \textcolor{gris}{\(F: U\ab \subseteq \R^n \to \R^m\) é \emph{diff} em \(p \in \R^n\) sse \(\exists L \in \hom(\R^n, \R^m): F(p+v) = F(p) + Lv + o\|v\|\), quando \(v\to 0\). }\\ 
(b) \(\blacksquare\) (\emph{TVM\(-\R^n\)}) \textcolor{gris}{Sejam \(f: U\ab \subseteq \R^n \to \R\) continua em \([p, p+v]\) e diff em \((p,p+tv)\), então \(\exists \theta \in (0,1)\) tal que \( f(p+v) - f(p) = df(p + \theta v) \cdot v \).   }\\ 
(c) \(\blacksquare\) (\emph{Função Inversa}) \textcolor{gris}{Sejam \(U\ab \subseteq \R^n\), \(F \in C^k(U,\R^n)\) e \(p\in U\). Então, \(^kF: \V_p \ \leftrightharpoons \ DF(p)\) é isomorfismo.  }\\ 
(d) Submersão. \textcolor{gris}{Uma aplicação \(F\) diff em \(p\) é submersão (em \(p\)) sse \(DF(p): \R^n \twoheadrightarrow \R^m\). }\\
(e) \(\blacksquare\) (\emph{FLS}) \textcolor{gris}{Seja \(F\in C^k(U,\R^m)\) submersão em \(p\in U\). Então, \(\exists \V_p \subseteq U\) e \(G:V\ab \simeq \V_p\) tais que \(F\circ G: V \ni (\x, x_{m+1}, \ldots, x_n) \mapsto \x \in \R^m\). }\\
(f) \(\blacksquare\) (\emph{Função Implícita}) \textcolor{gris}{Sejam \(U\ab = V\times W \subseteq \R^n \times \R^m\), \(G \in C^k(U,\R^m)\) e \((p,q) \in U\). Se \(G(p,q)=0\) e \(\det JG_q(p,q) \neq 0\), então \(\exists \mathcal{W}_q\) e \(\exists F \in C^k(\V_p, \R^m)\) tais que \(\forall (\x,\y) \in \V_p\times \mathcal{W}_q, \ G(\x,\y) = G(\x, F(\x)) = 0\). }\\ 
2. Sejam \(F,G: U\ab \subseteq \R^n \to \R^n\) diff em \(p \in U\), tais que \(F(p) = G(p)\). Mostre que,
\[DF(p) = DG(p) \ \leftrightharpoons \ \lim_{v\to 0} \frac{F(p+v) - G(p+v)}{\|v\|} = 0. \]
\textcolor{gris}{
    Segue essencialmente da definição, pois 
    \[DF(p) - DG(p) = \lim_{v\to 0} \frac{F(p+v) -DG(p+v)}{\|v\|}\] 
}
\hspace{-0.1cm}3. Sejam \(U:= \{(x,y) \in \R^2 : x,y >0 \}\), \(p,q >1\) tais que \(\frac{1}{p} + \frac{1}{q} = 1\) e \(f: U \to \R\) tal que, 
\[f(x,y) := \frac{x^p}{p} + \frac{y^q}{q}.\] 
(a) Explique por qué \(\forall t>0\), a restrição de \(f\) à hipérbole \(xy = t\) possui um ponto mínimo global \((x_t, y_t)\). \\ 
\textcolor{gris}{O conjunto \(\{(x,y) \in \R^2 : xy=t\}=: H_t\ce\subseteq \R^2 \) e ilimitado, mas \(f\) é contínua e \(f(x,y) \to \infty\) quando \(\|(x,y)\| \to \infty\), por tanto \(f\) tem mínimo global \((x_t, y_t) \in U\). } \\
(b) Utilizando Mult. de Lagrange, mostre que \(x_t^p = y_t^q\). \\ 
\textcolor{gris}{Seja \(g(x,y) = xy - t\), então \(\exists \lambda \in \R \) tal que,  %: df(x_t, y_t) = \lambda dg(x_t, y_t)\) e \(g(x_t,y_t) = 0\). 
\[ \begin{cases}
        x_t^{p-1}dx + y_t^{q-1}dy = \lambda (y_t dx + x_tdy) \\ 
        x_ty_t  - t = 0
    \end{cases}  \]
Isso \(\leftrightharpoons \ x_t^p = \lambda y_tx_t = \lambda x_ty_t = y_t^q = \lambda t  \).
}

(c) Mostre que \(f(x_t, y_t) = t \) e conclua que \(\forall x,y > 0\), 
\[ xy \leq \frac{x^p}{p} + \frac{y^q}{q}.  \]
\- \vspace{-0.5cm}
\textcolor{gris}{
\begin{align*}
f(x_t, y_t) &=  \frac{x_t^p}{p} + \frac{y_t^q}{q} = x_t^p \left(\frac{1}{p} + \frac{1}{q}\right) = x_t^p\\ 
&= (x_t^p)^{\frac{1}{p} + \frac{1}{q}} = x_t (x_t^{p})^{\frac{1}{q}} = x_ty_t = t 
\end{align*}
Dados \(x,y > 0\), \(xy = t = f(x_t, y_t) \leq f(x,y) = \frac{x^p}{p} + \frac{y^q}{q}\). 
}\\ 
4. Utilize \(\blacksquare\) (\emph{Função Implícita}) para mostrar que o sistema, 
\[\begin{cases}
    \textcolor{gray}{G_1(\x) = } \  w^2 + 2x^2+y^2 - z^2 - 6 = 0 \\ 
    \textcolor{gray}{G_1(\x) = } \  wxy - xyz = 0 
\end{cases}, \]
pode ser resolvido em termos de \(w = w(y,z)\) e \(x = x(y,z)\) numa vizinhança do ponto \((p,q) = ((y,z),(x,w)) = ((2,1),(1,1))\). Calcule as derivadas parciais de \(w\) e \(x\) em \(p\). \\
\textcolor{gris}{
    Sejam \(G(\x,\y) = (G_1(\x,\y), G_2(\x,\y)): \R^4 \to \R^2\). Note que \(G \in C^\infty (\R^4, \R^2) \), \(G(p,q) = (0,0)\) e,
    \[\det (JG_q(p,q)) = \det \begin{bmatrix}
        4 & 2 \\ 
        0 & 2
    \end{bmatrix} = 8 \neq 0.\]  
    Pelo \(\blacksquare\) (\emph{Função Implícita}) \(\exists \mathcal{W}_q\) e \(F\in C^\infty(\V_p, \R^2)\) onde \(\forall (\x, \y) \in \V_p \times \mathcal{W}_q, \ G(\x,\y ) = G(\x, F(\x))\), seja lá
    \[F(y,z) = (x(y,z), w(y,z)).\]  
    Pelo {\scriptsize \(\boxtimes\)} 8.5. \(JG_q(p,q)\cdot JF_p(p,q) = - JG_p(p,q)\), ou seja,
    \[ \begin{bmatrix}
        x_y(p,q) & x_z(p,q) \\
        w_y(p,q) & w_z(p,q) 
    \end{bmatrix}  = \begin{bmatrix}
        4 & 2 \\ 
        0 & 2
    \end{bmatrix}^{-1} \begin{bmatrix}
        -4 & 2 \\ 
        0 & 2
    \end{bmatrix} = \begin{bmatrix}
        -1 & 0 \\ 
        0& 1
    \end{bmatrix}.\]
}
5. Mostre que \(\omega \in \Omega^1(\R^2\setminus \{0\}) \) é fechada mas não exata,  
\[\omega := -\frac{y}{x^2 + y^2}dx + \frac{x}{x^2+y^2} dy \textcolor{gray}{ \ = fdx + gdy}.\] 
\textcolor{gris}{
    Lembre-se que \(\omega \) é \underline{fechada} \(\leftrightharpoons\ d\omega = 0 \), 
    \[ d\omega = (g_x - f_y) dx \wedge dy = \frac{2[x^2+y^2] - 2x^2 - 2y^2}{[x^2+y^2]^2}dx \wedge dy = 0. \] 
    Seja \(\E^1 \subseteq \R^2\), variedade compacta, orientada, sem borda, suponha por contradição que \(\omega \) for \underline{exata}, usando \(\psi: [0,2\pi) \ni t \mapsto (\cos(t), \sin(t)) \in \E^1\) temos,% então \(\int_{\E^1} \omega =0\), pelo {\scriptsize \(\boxtimes\)} (\emph{Stokes}) mas,  
    \[\int_{\E^1}w = \int_0^{2\pi} \sin^2(t) + \cos^2(t)  dt = \int_0^{2\pi}  dt = 2\pi \neq 0. \ \lightning\]
}
\hspace{-0.15cm}De um exemplo de \(_d\overline{M}\subseteq \R^n\) compacta que possui uma forma de volume exata. \\ 
\textcolor{gris}{
    Seja \(_2\overline{M} = \overline{\B^2} := \{\x \in \R^2: \|\x\|\leq 1\}\). Note que \(\omega_\text{vol} = dx \wedge dy\) e \(w = d(-\frac{y}{2}dx+\frac{x}{2}dy)\). 
}\\
Porém, se \(M\) for compacta e \(\partial M = \emptyset \), então nenhuma forma de volume é exata. \\
\textcolor{gris}{
    Supõe por contradição que \(\omega_\text{vol} = d\eta \), então do \(\blacksquare\) (\emph{Stokes}), 
    \[0< \int_M\omega_\text{vol} = \int_{\partial M} \eta = 0. \ \lightning\] 
}
